\section{Regularized modular rows and all derivative evaluations}
\label{sec:modular}

The modular draft obtains even-weight Eisenstein series by pairing
polygamma values along lattice rows. Its weight-two question requires
care: the paired rows converge rapidly, although the ungrouped lattice
sum is not absolutely convergent. We first prove an explicit bound for
every differentiated row. We then determine the change of summation
order directly, before using any modular transformation law. Finally,
Ramanujan's differential system evaluates all the resulting derivative
rows at the square and hexagonal CM points.

Throughout this section, let
\[
\tau\in\mathbb H:=\{z\in\C:\Impart z>0\},\qquad
q=e^{2\pi\ii\tau},\qquad r=|q|=e^{-2\pi\Impart\tau}<1,
\qquad \rho=\tau_{\mathrm{hex}}=\frac{1+\ii\sqrt3}{2}.
\]
For even \(w\ge2\) and an integer \(j\ge0\), put
\begin{equation}
\mathcal T_{w,j}(\tau)
=\sum_{n\ge1}n^j\left[
 \psi^{(w-1+j)}(n\tau)
 +(-1)^j\psi^{(w-1+j)}(1-n\tau)\right].
\label{eq:modular-T-definition}
\end{equation}
Convergence will follow from the next theorem. The sign on the reflected
term is part of the definition; differentiating that term changes its
sign at every step.

\subsection{Rational row formulas and explicit tail bounds}

The negative-index polylogarithms appearing below are elementary
rational functions. For \(d\ge0\), define
\[
R_d(t)=\Li_{-d}(t),\qquad
R_0(t)=\frac{t}{1-t},\qquad
R_{d+1}(t)=t\frac{d}{dt}R_d(t).
\]
These relations are identities of rational functions, since on
\(|t|<1\) they follow by differentiating
\(R_d(t)=\sum_{\ell\ge1}\ell^d t^\ell\). In particular,
\[
R_1(t)=\frac{t}{(1-t)^2},\qquad
R_2(t)=\frac{t(1+t)}{(1-t)^3},\qquad
R_3(t)=\frac{t(1+4t+t^2)}{(1-t)^4}.
\]

\begin{theorem}[Rational rows and an explicit remainder]
\label{thm:modular-rows-tail}
For \(w,j,\tau\) as above,
\begin{align}
&\psi^{(w-1+j)}(n\tau)
 +(-1)^j\psi^{(w-1+j)}(1-n\tau)
 =(2\pi\ii)^{w+j}\Li_{1-w-j}(q^n),
\label{eq:modular-single-row}\\
&\mathcal T_{w,j}(\tau)
 =(2\pi\ii)^{w+j}\sum_{n\ge1}n^j\Li_{1-w-j}(q^n).
\label{eq:modular-rational-rows}
\end{align}
The series converges normally on every compact subset of \(\mathbb H\).
If \(\mathcal T_{w,j}^{[N]}\) is its sum through \(n=N\), where
\(N\ge0\), then
\begin{equation}
\left|\mathcal T_{w,j}(\tau)-\mathcal T_{w,j}^{[N]}(\tau)\right|
\le (2\pi)^{w+j}(N+1)^j
 \frac{\Li_{-j}(r)}{r}\,
 \Li_{1-w-j}\!\left(r^{N+1}\right).
\label{eq:modular-tail-bound}
\end{equation}
Thus the bound contains only rational functions of \(r\) and
\(r^{N+1}\), together with an explicit power of \(\pi\).
\end{theorem}

\begin{proof}
For \(\Impart z>0\), splitting the absolutely convergent bilateral
power sum into its nonnegative and negative indices gives
\begin{equation}
\sum_{m\in\Z}(m+z)^{-w}
=\frac{\psi^{(w-1)}(z)+\psi^{(w-1)}(1-z)}{(w-1)!}.
\label{eq:modular-bilateral-polygamma}
\end{equation}
Here evenness of \(w\) is essential. The classical cotangent expansion
and its partial-fraction interpretation~\cite{DLMF} read
\[
\pi\cot(\pi z)=-\pi\ii
 -2\pi\ii\sum_{\ell\ge1}e^{2\pi\ii\ell z},\qquad
\pi\cot(\pi z)=\lim_{M\to\infty}\sum_{m=-M}^{M}\frac1{m+z}.
\]
Differentiate \(w-1\) times. On compact subsets of the upper
half-plane the exponential series and its derivatives converge
uniformly. After the first derivative, the bilateral series and all
its subsequent derivatives also converge normally. Since
\((-1)^{w-1}=-1\), the two minus signs cancel and yield
\[
\psi^{(w-1)}(z)+\psi^{(w-1)}(1-z)
=(2\pi\ii)^w\sum_{\ell\ge1}\ell^{w-1}e^{2\pi\ii\ell z}.
\]
Differentiating another \(j\) times with respect to \(z\) proves
\eqref{eq:modular-single-row}. Multiplication by \(n^j\) gives the
summand in \eqref{eq:modular-rational-rows}.

To prove the remainder bound, expand the rational polylogarithm in
its nonnegative majorizing series. The absolute value of the omitted
rows is at most
\begin{equation}
(2\pi)^{w+j}\sum_{\ell\ge1}\ell^{w+j-1}
                  \sum_{n>N}n^j r^{n\ell}.
\label{eq:modular-positive-majorant}
\end{equation}
Write \(n=N+1+a\). For \(a\ge0\) and \(\ell\ge1\),
\[
n^j\le(N+1)^j(a+1)^j,\qquad r^{a\ell}\le r^a.
\]
Consequently
\[
\sum_{n>N}n^j r^{n\ell}
\le (N+1)^j r^{(N+1)\ell}
       \sum_{a\ge0}(a+1)^j r^a
=(N+1)^j r^{(N+1)\ell}\frac{\Li_{-j}(r)}{r}.
\]
Substitution in \eqref{eq:modular-positive-majorant} proves
\eqref{eq:modular-tail-bound}. In particular, the majorant is finite.
On a compact subset of \(\mathbb H\), one has \(r\le r_0<1\), and
each of its power series with nonnegative coefficients is bounded by
its value at \(r_0\). This proves normal convergence, for every fixed
\(w,j\), and justifies all subsequent termwise differentiations.
\end{proof}

For \(j=0\), the tail formula simplifies to
\[
\left|\mathcal T_{w,0}-\mathcal T_{w,0}^{[N]}\right|
\le\frac{(2\pi)^w}{1-r}\Li_{1-w}\!\left(r^{N+1}\right).
\]
The case \(w=2\), which will be particularly useful, is
\begin{equation}
\left|\mathcal T_{2,0}-\mathcal T_{2,0}^{[N]}\right|
\le\frac{4\pi^2r^{N+1}}
 {(1-r)(1-r^{N+1})^2}.
\label{eq:modular-weight-two-tail}
\end{equation}
Theorem~\ref{thm:modular-rows-tail} also proves
\begin{equation}
\frac{d^j}{d\tau^j}\mathcal T_{w,0}(\tau)
=\mathcal T_{w,j}(\tau).
\label{eq:modular-termwise-derivatives}
\end{equation}
Computationally, \eqref{eq:modular-single-row} is preferable to
evaluating its two polygamma terms separately. Their algebraic
asymptotic contributions cancel, leaving the exponentially small
paired row. The rational formula performs that cancellation exactly
before numerical evaluation.

\subsection{The weight-two ordering correction from telescoping}

For \(w\ge4\), the ordinary absolutely convergent lattice series is
\[
G_w(\tau)=\sum_{(m,n)\in\Z^2\setminus\{(0,0)\}}
                  (m+n\tau)^{-w}
=2\zeta(w)+\frac{2}{(w-1)!}\mathcal T_{w,0}(\tau).
\]
For \(w=2\), specify the order explicitly:
\begin{align}
G_2(\tau)&=\sum_{n\in\Z}
        \left(\sum_{m\in\Z}'(m+n\tau)^{-2}\right)
 =\frac{\pi^2}{3}+2\mathcal T_{2,0}(\tau),
\label{eq:modular-G2-definition}\\
C_2(\tau)&=\sum_{m\in\Z}
        \left(\sum_{n\in\Z}'(m+n\tau)^{-2}\right).
\label{eq:modular-C2-definition}
\end{align}
The prime omits only \((m,n)=(0,0)\). In the first expression, each
inner row and the outer series of completed rows converge absolutely.
The reverse expression exists as well: changing the names and signs
of the one-dimensional indices gives
\begin{equation}
C_2(\tau)=\tau^{-2}G_2(-1/\tau).
\label{eq:modular-C2-coordinate}
\end{equation}
This is a coordinate identity for the prescribed iterated sums, not
an assumption of modularity. The series of absolute values over all
lattice points diverges, so an interchange of the two orders still
requires a calculation.

\begin{theorem}[Summation-order anomaly]
\label{thm:modular-ordering}
For every \(\tau\in\mathbb H\),
\begin{equation}
C_2(\tau)=G_2(\tau)-\frac{2\pi\ii}{\tau}.
\label{eq:modular-anomaly}
\end{equation}
Consequently
\begin{equation}
G_2(-1/\tau)=\tau^2G_2(\tau)-2\pi\ii\tau,
\qquad G_2(\tau+1)=G_2(\tau).
\label{eq:modular-G2-transform}
\end{equation}
\end{theorem}

\begin{proof}
We prove the difference of the two orders by replacing their common
part with an absolutely convergent series. Put
\[
h(z)=\frac1z-\frac1{z+1},\qquad
f(z)=\frac1{z^2}-h(z)=\frac1{z^2(z+1)},
\]
and define
\[
F(\tau)=\sum_{\substack{n\in\Z\setminus\{0\}\\m\in\Z}}
                    f(m+n\tau).
\]
No pole occurs in this sum. The real-linear map
\((m,n)\mapsto m+n\tau\) is invertible, so for some \(c_\tau>0\),
\(|m+n\tau|\ge c_\tau\sqrt{m^2+n^2}\).
Outside a finite set, \(|f(z)|\le2|z|^{-3}\); hence \(F\) converges
absolutely. It can therefore be summed in either order.

For each fixed \(n\ne0\), the series \(\sum_m h(m+n\tau)\) is
absolutely convergent and its symmetric partial sum telescopes:
\[
\sum_{m=-M}^{M}h(m+n\tau)
=\frac1{-M+n\tau}-\frac1{M+1+n\tau}\longrightarrow0.
\]
Separating the row \(n=0\) proves
\begin{equation}
G_2(\tau)=\frac{\pi^2}{3}+F(\tau).
\label{eq:modular-G2-absolute-part}
\end{equation}

For the reverse order, define for integer \(m\)
\[
A(m)=\lim_{K\to\infty}
              \sum_{0<|n|\le K}\frac1{m+n\tau}.
\]
This symmetric limit exists: pairing \(n\) and \(-n\) gives an
\(O(n^{-2})\) summand. The cotangent partial-fraction formula gives
\begin{equation}
A(m)=\frac{\pi}{\tau}\cot\!\left(\frac{\pi m}{\tau}\right)-\frac1m
\quad(m\ne0),\qquad A(0)=0.
\label{eq:modular-A-cotangent}
\end{equation}
The value at zero is also the removable limit of the expression on
the right. For each fixed \(m\), the \(n\)-series of \(h(m+n\tau)\)
is absolutely convergent. Taking its symmetric partial sums shows
\[
\sum_{n\ne0}h(m+n\tau)=A(m)-A(m+1).
\]
It follows that
\begin{align*}
C_2(\tau)-\frac{\pi^2}{3}-F(\tau)
&=\lim_{M\to\infty}\sum_{m=-M}^{M}\sum_{n\ne0}h(m+n\tau)\\
&=\lim_{M\to\infty}\bigl(A(-M)-A(M+1)\bigr).
\end{align*}
The extraction of \(F\) here is legitimate by its absolute
convergence; no rearrangement of the original double power sum has
been made.

Since \(\Impart(1/\tau)<0\),
\[
\cot(\pi m/\tau)\longrightarrow\ii\quad(m\to+\infty),
\qquad
\cot(\pi m/\tau)\longrightarrow-\ii\quad(m\to-\infty).
\]
Equation~\eqref{eq:modular-A-cotangent} therefore gives
\(A(-M)-A(M+1)\to-2\pi\ii/\tau\). Combining this with
\eqref{eq:modular-G2-absolute-part} proves
\eqref{eq:modular-anomaly}. Equation~\eqref{eq:modular-C2-coordinate}
now gives the first transformation in
\eqref{eq:modular-G2-transform}. Periodicity follows directly from
the \(q\)-series in \eqref{eq:modular-rational-rows}.
\end{proof}

The anomaly is classical; Romik and Scherer discuss it together with
more general shape-dependent summation procedures~\cite{RomikScherer}.
The proof above fixes both its sign and the normalization needed for
the draft's polygamma sums.

\begin{corollary}[The missing weight-two CM rows]
\label{cor:modular-weight-two-CM}
One has
\begin{align}
G_2(\ii)&=\pi,&G_2(\rho)&=\frac{2\pi}{\sqrt3},
\label{eq:modular-G2-CM}\\
\mathcal T_{2,0}(\ii)&=\frac\pi2-\frac{\pi^2}{6},&
\mathcal T_{2,0}(\rho)&=\frac\pi{\sqrt3}-\frac{\pi^2}{6}.
\label{eq:modular-trigamma-CM}
\end{align}
\end{corollary}
\begin{proof}
At \(\tau=\ii\), the first transformation in
\eqref{eq:modular-G2-transform} gives
\(G_2(\ii)=-G_2(\ii)+2\pi\).
For \(\rho\), use \(-1/\rho=\rho-1\) and periodicity. The same
transformation becomes
\[
(\rho^2-1)G_2(\rho)=2\pi\ii\rho.
\]
Since \(\rho^2-1=\ii\sqrt3\rho\), this proves the second value in
\eqref{eq:modular-G2-CM}. Substitution into
\eqref{eq:modular-G2-definition} proves
\eqref{eq:modular-trigamma-CM}.
\end{proof}

\subsection{CM seeds and the complete differential tower}

Normalize the Eisenstein series by
\[
E_2=\frac{3G_2}{\pi^2},\qquad
E_4=\frac{G_4}{2\zeta(4)},\qquad
E_6=\frac{G_6}{2\zeta(6)}.
\]
The seed values, in precisely this normalization, are
\begin{align}
(E_2(\ii),E_4(\ii),E_6(\ii))
&=\left(\frac3\pi,
 \frac{3\Gamma(1/4)^8}{64\pi^6},0\right),
\label{eq:modular-square-seeds}\\
(E_2(\rho),E_4(\rho),E_6(\rho))
&=\left(\frac{2\sqrt3}\pi,0,
 \frac{27\Gamma(1/3)^{18}}{512\pi^{12}}\right).
\label{eq:modular-hex-seeds}
\end{align}
We have already proved the weight-two entries. The zero entries follow
by rotating the absolutely convergent lattices. Multiplication by
\(\ii\) fixes the square lattice and sends its weight-six sum to its
negative; multiplication by \(\rho\) fixes the hexagonal lattice and
multiplies its weight-four sum by \(\rho^{-4}\ne1\).

For completeness, the classical period integrals fix the nonzero
entries without numerical identification. The standard inverse
Weierstrass integral~\cite{DLMF}, with
\(g_2=60G_4\) and \(g_3=140G_6\), gives positive real periods
\begin{align*}
p_{\mathrm{sq}}
&=2\int_1^\infty\frac{dx}{\sqrt{4x^3-4x}}
 =2\int_0^1\frac{du}{\sqrt{1-u^4}}
 =\frac12 B\!\left(\frac14,\frac12\right)
 =\frac{\Gamma(1/4)^2}{2\sqrt{2\pi}},\\
p_{\mathrm{hex}}
&=2\int_1^\infty\frac{dx}{\sqrt{4x^3-4}}
 =\frac13 B\!\left(\frac16,\frac12\right)
 =\frac{\Gamma(1/3)^3}{2^{4/3}\pi}.
\end{align*}
The corresponding lattices are
\(p_{\mathrm{sq}}(\Z+\ii\Z)\) for \((g_2,g_3)=(4,0)\), and
\(p_{\mathrm{hex}}(\Z+\rho\Z)\) for \((g_2,g_3)=(0,4)\).
Their square and equianharmonic shapes follow from the rotations of
the respective cubics; the real period integral fixes the scale and
orientation. The first substitution above is \(x=u^{-2}\); in the
second beta integral use \(t=x^{-3}\). Euler's reflection and
duplication formulas give the displayed gamma evaluations.
Scaling the lattice invariants consequently yields
\[
G_4(\ii)=\frac{p_{\mathrm{sq}}^4}{15}
 =\frac{\Gamma(1/4)^8}{960\pi^2},\qquad
G_6(\rho)=\frac{p_{\mathrm{hex}}^6}{35}
 =\frac{\Gamma(1/3)^{18}}{8960\pi^6}.
\]
Using \(2\zeta(4)=\pi^4/45\) and
\(2\zeta(6)=2\pi^6/945\) proves the remaining entries of
\eqref{eq:modular-square-seeds}--\eqref{eq:modular-hex-seeds}.

Ramanujan's classical differential identities~\cite{Ramanujan} are
\begin{equation}
\theta E_2=\frac{E_2^2-E_4}{12},\qquad
\theta E_4=\frac{E_2E_4-E_6}{3},\qquad
\theta E_6=\frac{E_2E_6-E_4^2}{2},\qquad
\theta=\frac1{2\pi\ii}\frac{d}{d\tau}.
\label{eq:modular-Ramanujan-system}
\end{equation}
Encode these identities by the rational derivation
\begin{equation}
\mathcal V
=\frac{X^2-Y}{12}\frac{\partial}{\partial X}
 +\frac{XY-Z}{3}\frac{\partial}{\partial Y}
 +\frac{XZ-Y^2}{2}\frac{\partial}{\partial Z},
\qquad P_0=X,\quad P_{j+1}=\mathcal V P_j.
\label{eq:modular-P-recurrence}
\end{equation}
For example,
\begin{align*}
P_1&=\frac{X^2-Y}{12},\\
P_2&=\frac{X^3-3XY+2Z}{72},\\
P_3&=\frac{X^4-6X^2Y+8XZ-3Y^2}{288}.
\end{align*}

\begin{theorem}[All differentiated weight-two rows]
\label{thm:modular-all-jets}
Set \(\mathcal S_j(\tau)=\mathcal T_{2,j}(\tau)\). For every
\(\tau\in\mathbb H\),
\begin{align}
\mathcal S_0(\tau)&=\frac{\pi^2}{6}\bigl(E_2(\tau)-1\bigr),
\label{eq:modular-S0}\\
\mathcal S_j(\tau)&=\frac{\pi^2}{6}(2\pi\ii)^j
 P_j(E_2(\tau),E_4(\tau),E_6(\tau)),\qquad j\ge1.
\label{eq:modular-all-jets}
\end{align}
In particular, substitution of
\eqref{eq:modular-square-seeds}--\eqref{eq:modular-hex-seeds}
evaluates every \(\mathcal S_j(\ii)\) and
\(\mathcal S_j(\rho)\) by a finite polynomial calculation over
\(\Q\) and the indicated gamma periods.
\end{theorem}

\begin{proof}
Equation~\eqref{eq:modular-G2-definition} and the definition of
\(E_2\) give \eqref{eq:modular-S0}. By the chain rule and
\eqref{eq:modular-Ramanujan-system}, every polynomial
\(P\in\Q[X,Y,Z]\) satisfies
\[
\theta\bigl(P(E_2,E_4,E_6)\bigr)
=(\mathcal V P)(E_2,E_4,E_6).
\]
Induction from \(P_0=X\) therefore gives
\(\theta^jE_2=P_j(E_2,E_4,E_6)\).
Differentiate \eqref{eq:modular-S0} \(j\ge1\) times and use
\eqref{eq:modular-termwise-derivatives}. The constant term disappears,
and \(d/d\tau=2\pi\ii\theta\) contributes the factor
\((2\pi\ii)^j\). This proves \eqref{eq:modular-all-jets}; the CM
specializations use the established seed values.
\end{proof}

For clarity, the first two differentiated evaluations are
\begin{align}
\sum_{n\ge1}n\left[\psi''(n\ii)-\psi''(1-n\ii)\right]
 &=\frac{\ii\pi}{4}
   -\frac{\ii\Gamma(1/4)^8}{768\pi^3},
\label{eq:modular-S1-square}\\
\sum_{n\ge1}n\left[\psi''(n\rho)-\psi''(1-n\rho)\right]
 &=\frac{\ii\pi}{3},
\label{eq:modular-S1-hex}\\
\sum_{n\ge1}n^2\left[\psi'''(n\ii)+\psi'''(1-n\ii)\right]
 &=-\frac\pi4+\frac{\Gamma(1/4)^8}{256\pi^3},
\label{eq:modular-S2-square}\\
\sum_{n\ge1}n^2\left[\psi'''(n\rho)+\psi'''(1-n\rho)\right]
 &=-\frac{2\sqrt3\pi}{9}
   -\frac{\Gamma(1/3)^{18}}{1024\pi^8}.
\label{eq:modular-S2-hex}
\end{align}
These follow by substituting the seed values into \(P_1,P_2\), so
their derivation contains no integer-relation search.

\begin{corollary}[Polynomial degree bounds for the CM row formulas]
\label{cor:modular-CM-degree}
Put
\(A=\Gamma(1/4)^8/\pi^4\) and
\(B=\Gamma(1/3)^{18}/\pi^9\). For each \(j\ge1\), there are
polynomials \(U_j\in\Q[t]\) and
\(V_j\in\Q(\sqrt3)[t]\) such that
\[
\frac{\mathcal S_j(\ii)}{\ii^j\pi}=U_j(A),\qquad
\frac{\mathcal S_j(\rho)}{\ii^j\pi}=V_j(B),
\]
with
\[
\deg U_j\le\left\lfloor\frac{j+1}{2}\right\rfloor,
\qquad
\deg V_j\le\left\lfloor\frac{j+1}{3}\right\rfloor.
\]
\end{corollary}
\begin{proof}
Give \(X,Y,Z\) weights \(2,4,6\). Each term of
\(\mathcal V\) raises weight by two, so \(P_j\) is homogeneous
of weight \(2j+2\). At the square point, the surviving monomials
\(X^aY^b\) have \(a+2b=j+1\). On inserting the seeds into
\eqref{eq:modular-all-jets} and dividing by \(\ii^j\pi\), each
becomes a rational multiple of \(A^b\). At the hexagonal point,
the surviving monomials \(X^aZ^c\) satisfy \(a+3c=j+1\), and
the same substitution produces an element of
\(\Q(\sqrt3)B^c\). The stated degree bounds follow.
\end{proof}

\subsection{What this resolves and how it is verified}

Corollary~\ref{cor:modular-weight-two-CM} resolves the weight-two
row-sum question posed in the modular draft. Theorem~\ref{thm:modular-all-jets}
extends it to all polynomially weighted derivative rows, with the
explicit error estimate of Theorem~\ref{thm:modular-rows-tail}.
These are evaluations of specified infinite aggregates. They do not
evaluate an isolated component such as \(\Impart\psi'(\ii)\), nor
do they prove that such a component is an Eichler integral. Any proposed
identification of that isolated constant needs its own exact transform
or functional equation.

The accompanying \path{code/modular_rows.py} constructs the rational
functions \(R_d\) and the polynomials \(P_j\) by exact arithmetic.
It checks individual rows against a separate complex-polygamma
implementation, compares the CM formulas for \(j=0,\ldots,4\) with
their rapidly convergent row sums, and checks representative tails.
The output \path{data/modular_rows.json} records the exact coefficients,
analytic truncation bounds, and numerical residuals. Its decimal
calculations use ordinary high-precision arithmetic; they are numerical
checks, not outward-rounded interval certificates. A machine enclosure
can be obtained by evaluating the rational formulas and
\eqref{eq:modular-tail-bound} with directed interval arithmetic. The
identities themselves are proved above and do not depend on the
numerical checks.
